\section{Main Findings}
\label{sec:results}
\subsection{RQ1 Finding: Process Savings Do Not Establish Extra Capacity}
Stage-0 PSS decreases by 382.56, 382.56 and 382.55\,MiB in the three release
cycles, with same-precision transfer checks passing. At the same observations,
node \texttt{MemAvailable} changes by $-5.82$, $-1.12$ and $-11.40$\,MiB.
The separate artifact appendix plots both endpoints. Unmapping works, but cached
file pages and other node activity are not causally separated. Additional work
admitted by the saving is not demonstrated.


The state study also distinguishes successful transfer from target-model accuracy. Retention has
mean document KL 0.0835 against Q4, while rebuilding recovers Q4 logits at the
measured positions. Complete evaluation takes 25.70 versus 68.16 minutes, a
2.65$\times$ ratio. This includes separate teacher-forced and generated passes,
not one application's request latency. Two answers are clearly supported across
policies. A third shares a grounding problem. Mean F1 equality is not semantic
equivalence, and the three-document pilot cannot establish general preservation.

The MiniLM controls point to a different bottleneck. The real 81-chunk update
publishes in 3.64\,s alone versus 11.77\,s concurrently with Q4 in the earlier
decode-overlap control. Embedding explains approximately 97--99\% of publication
time. The update progresses without a failure caused by insufficient capacity. CPU scheduling,
cache and bandwidth effects remain incompletely separated. This supports
investigating execution interference, not attributing the slowdown to a useful
memory-release opportunity. These timings are not pooled with the later
prefill-overlap application experiment.

\subsection{RQ2 Findings: Backend Settings and CPU Allocation}
The prototype tests select repacked microbatch 64. Its complete
fresh-prefix ladder is approximately 2.0--2.1$\times$ faster than raw microbatch
2 from 128 to 1024 tokens (Figure~\ref{fig:research-performance}a). The underlying
mechanisms already exist in the runtime. In the later comparison, all four
standard candidates pass the old supported fixtures. The selected standard
configuration is competitive with both private prototypes
(Figure~\ref{fig:research-performance}b).
\begin{figure*}[t]
\centering
\begin{tikzpicture}
\begin{axis}[name=left,width=.47\textwidth,height=4.5cm,xlabel={Prompt tokens},ylabel={Prefill compute (s)},
 xmin=100,xmax=1060,ymin=0,ymax=600,xtick={128,256,512,1024},ytick={0,200,400,600},
 axis lines=left,grid=major,grid style={black!8},tick label style={font=\scriptsize},label style={font=\small},
 title={(a) Opportunity against a restricted baseline},title style={font=\small},
 legend style={font=\scriptsize,draw=none,at={(.03,.97)},anchor=north west}]
\addplot[black!60,dashed,mark=square*,line width=1pt] table[x=tokens,y=raw]{generated/ladder.dat};
\addlegendentry{Raw microbatch 2}
\addplot[answerblue,mark=*,line width=1.2pt] table[x=tokens,y=packed]{generated/ladder.dat};
\addlegendentry{Repacked microbatch 64}
\end{axis}
\begin{axis}[at={(left.east)},anchor=west,xshift=1.12cm,width=.47\textwidth,height=4.5cm,
 ybar,bar width=10pt,xmin=-.55,xmax=2.55,ymin=0,ymax=205,xtick={0,1,2},
 xticklabels={Standard,Full,Skip4},ylabel={Complete answer (s)},
 axis lines=left,ymajorgrids,grid style={black!8},tick label style={font=\scriptsize},label style={font=\small},
 title={(b) Comparison with a tuned standard server},title style={font=\small},
 legend style={font=\scriptsize,draw=none,at={(.5,.98)},anchor=north,legend columns=2},
 nodes near coords,nodes near coords style={font=\scriptsize},point meta=y]
\addplot[fill=passgreen!80,draw=passgreen] table[x=id,y=short]{generated/latency.dat};\addlegendentry{184 tokens}
\addplot[fill=answerblue!80,draw=answerblue] table[x=id,y=long]{generated/latency.dat};\addlegendentry{565 tokens}
\end{axis}
\end{tikzpicture}
\caption{The baseline changes the research interpretation. (a) Complete three-run
fresh-KV prefix medians show approximately 2.0--2.1$\times$ prefill acceleration
with existing CPU mechanisms. Synthetic prefixes measure execution, not answer
quality. (b) In the later matched comparison, a tuned unmodified server matches
or exceeds the private full/skip4 prototypes on warm answer time. Each bar is a
three-run median for a fixed prompt, with no encoder overlap; startup is excluded.
Panels use their own archived campaigns and are not a pooled speedup comparison.}
\label{fig:research-performance}
\end{figure*}


\begin{table}[t]
\caption{Matched warm complete-answer medians (seconds), three runs per fixed
prompt, no encoder overlap. Standard is selected only on old fixtures.}
\label{tab:standard}
\centering
\begin{tabular}{lrr}
\toprule
Configuration & Short & Long\\
\midrule
Standard & 70.2 & 164.0 \\
Full prototype & 72.9 & 167.1 \\
Skip4 prototype & 74.3 & 169.6 \\
\bottomrule

\end{tabular}
\end{table}
Table~\ref{tab:standard} shows that the prototype's large advantage over the
restricted configuration is not an advantage over this standard backend.
The added skip4 variant is slightly slower in these cells. The result does not
support a stable attention-skip speedup over standard. Three observations do
not establish significance or equivalence. The comparison with the restricted configuration is therefore insufficient
to claim a new execution method.

Warm acceleration also changes resource demand. Across observed sessions,
raw-mb2 and repack-mb64 median readiness latencies are approximately 9 and 21\,s.
Post-load anonymous RSS reaches approximately 0.8 and 8.6\,GiB. The faster warm execution therefore comes with higher startup and memory costs. Shared file ownership
and practical headroom are not normalized, and no energy benefit is measured.

\begin{figure}[t]
\centering
\begin{tikzpicture}
\begin{axis}[width=\columnwidth,height=4.6cm,ybar,bar width=11pt,
 symbolic x coords={Prefill,Decode,Update},xtick=data,ymin=0,ymax=255,
 ylabel={Seconds},tick label style={font=\scriptsize},label style={font=\small},
 legend style={font=\scriptsize,at={(.5,1.02)},anchor=south,legend columns=2},
 enlarge x limits=.22,nodes near coords,
 nodes near coords style={font=\tiny,/pgf/number format/fixed,/pgf/number format/fixed zerofill,/pgf/number format/precision=1}]
\addplot[fill=answerblue!65,nodes near coords style={xshift=-4pt}] table[x=phase,y=shared]{generated/budget-phases.dat};
\addplot[fill=updatepurple!65,nodes near coords style={xshift=4pt}] table[x=phase,y=partition]{generated/budget-phases.dat};
\legend{Shared,Partition 2+2}
\end{axis}
\end{tikzpicture}
\caption{E1 long-prompt phase medians, three runs per policy. Reserving cores
reduces update time while increasing prefill time; generation changes little.
Each bar is the median of one phase across three runs.}
\label{fig:budget-phases}
\end{figure}

\textbf{Ingestion improves at the expense of prefill.} For E1's long prompt,
shared execution gives median answer/update times of 174.3/13.8\,s. With 2+2,
they become 246.4/7.0\,s. Prefill rises from 143.7 to 216.3\,s, while generation
stays near 31\,s (Figure~\ref{fig:budget-phases}). Ingestion alone takes 4.8\,s.
The partition reduces ingestion interference but slows inference.


\subsection{RQ3 Finding: Aggregate Scores Hide Different Failures}
\begin{table}[t]
\caption{New-article results: official mean F1 and frozen automatic task passes.
There are 12 answerable (A) and 12 unanswerable (U) cases per arm.}
\label{tab:new-quality}
\centering
\begin{tabular}{lrrr}
\toprule
Configuration & F1 & A pass & U pass\\
\midrule
Standard & 0.491 & 10/12 & 0/12 \\
Full prototype & 0.491 & 10/12 & 0/12 \\
Skip4 prototype & 0.634 & 8/12 & 2/12 \\
\bottomrule

\end{tabular}
\end{table}
Standard and full prototype each pass ten of the 24 automatic contracts and
have identical aggregate F1. Skip4 also passes ten but changes which cases pass
(Figure~\ref{fig:research-quality}). Its higher F1 is not proof of preserved
individual answers. The frozen paired gate rejects skip4 for losses on cases
08 and 16, both passed by standard.
\begin{figure}[t]
\centering
\begin{tikzpicture}[font=\scriptsize]
\node at (1.37,.67) {12 answerable articles};\node at (4.20,.67) {12 unanswerable articles};
\foreach \i in {0,...,23} {\pgfmathtruncatemacro{\n}{mod(\i,12)+1}\node[font=\tiny] at (\i*.235+.104,.40) {\n};}
\node[anchor=east] at (-.12,.13) {Standard};
\node[anchor=east] at (-.12,-.27) {Full};
\node[anchor=east] at (-.12,-.67) {Skip4};
\fill[passgreen] (0.000,0.000) rectangle ++(.208,.27);
\fill[passgreen] (0.235,0.000) rectangle ++(.208,.27);
\fill[passgreen] (0.470,0.000) rectangle ++(.208,.27);
\fill[passgreen] (0.705,0.000) rectangle ++(.208,.27);
\fill[passgreen] (0.940,0.000) rectangle ++(.208,.27);
\fill[passgreen] (1.175,0.000) rectangle ++(.208,.27);
\fill[passgreen] (1.410,0.000) rectangle ++(.208,.27);
\fill[failgray] (1.645,0.000) rectangle ++(.208,.27);
\fill[passgreen] (1.880,0.000) rectangle ++(.208,.27);
\fill[failgray] (2.115,0.000) rectangle ++(.208,.27);
\fill[passgreen] (2.350,0.000) rectangle ++(.208,.27);
\fill[passgreen] (2.585,0.000) rectangle ++(.208,.27);
\fill[formatamber] (2.820,0.000) rectangle ++(.208,.27);
\fill[formatamber] (3.055,0.000) rectangle ++(.208,.27);
\fill[formatamber] (3.290,0.000) rectangle ++(.208,.27);
\fill[failgray] (3.525,0.000) rectangle ++(.208,.27);
\fill[formatamber] (3.760,0.000) rectangle ++(.208,.27);
\fill[formatamber] (3.995,0.000) rectangle ++(.208,.27);
\fill[failgray] (4.230,0.000) rectangle ++(.208,.27);
\fill[formatamber] (4.465,0.000) rectangle ++(.208,.27);
\fill[formatamber] (4.700,0.000) rectangle ++(.208,.27);
\fill[formatamber] (4.935,0.000) rectangle ++(.208,.27);
\fill[formatamber] (5.170,0.000) rectangle ++(.208,.27);
\fill[formatamber] (5.405,0.000) rectangle ++(.208,.27);
\fill[passgreen] (0.000,-0.400) rectangle ++(.208,.27);
\fill[passgreen] (0.235,-0.400) rectangle ++(.208,.27);
\fill[passgreen] (0.470,-0.400) rectangle ++(.208,.27);
\fill[passgreen] (0.705,-0.400) rectangle ++(.208,.27);
\fill[passgreen] (0.940,-0.400) rectangle ++(.208,.27);
\fill[passgreen] (1.175,-0.400) rectangle ++(.208,.27);
\fill[passgreen] (1.410,-0.400) rectangle ++(.208,.27);
\fill[failgray] (1.645,-0.400) rectangle ++(.208,.27);
\fill[passgreen] (1.880,-0.400) rectangle ++(.208,.27);
\fill[failgray] (2.115,-0.400) rectangle ++(.208,.27);
\fill[passgreen] (2.350,-0.400) rectangle ++(.208,.27);
\fill[passgreen] (2.585,-0.400) rectangle ++(.208,.27);
\fill[formatamber] (2.820,-0.400) rectangle ++(.208,.27);
\fill[formatamber] (3.055,-0.400) rectangle ++(.208,.27);
\fill[formatamber] (3.290,-0.400) rectangle ++(.208,.27);
\fill[failgray] (3.525,-0.400) rectangle ++(.208,.27);
\fill[formatamber] (3.760,-0.400) rectangle ++(.208,.27);
\fill[formatamber] (3.995,-0.400) rectangle ++(.208,.27);
\fill[failgray] (4.230,-0.400) rectangle ++(.208,.27);
\fill[formatamber] (4.465,-0.400) rectangle ++(.208,.27);
\fill[formatamber] (4.700,-0.400) rectangle ++(.208,.27);
\fill[formatamber] (4.935,-0.400) rectangle ++(.208,.27);
\fill[formatamber] (5.170,-0.400) rectangle ++(.208,.27);
\fill[formatamber] (5.405,-0.400) rectangle ++(.208,.27);
\fill[passgreen] (0.000,-0.800) rectangle ++(.208,.27);
\fill[passgreen] (0.235,-0.800) rectangle ++(.208,.27);
\fill[passgreen] (0.470,-0.800) rectangle ++(.208,.27);
\fill[passgreen] (0.705,-0.800) rectangle ++(.208,.27);
\fill[failgray] (0.940,-0.800) rectangle ++(.208,.27);
\fill[passgreen] (1.175,-0.800) rectangle ++(.208,.27);
\fill[passgreen] (1.410,-0.800) rectangle ++(.208,.27);
\fill[failgray] (1.645,-0.800) rectangle ++(.208,.27);
\fill[failgray] (1.880,-0.800) rectangle ++(.208,.27);
\fill[failgray] (2.115,-0.800) rectangle ++(.208,.27);
\fill[passgreen] (2.350,-0.800) rectangle ++(.208,.27);
\fill[passgreen] (2.585,-0.800) rectangle ++(.208,.27);
\fill[formatamber] (2.820,-0.800) rectangle ++(.208,.27);
\fill[formatamber] (3.055,-0.800) rectangle ++(.208,.27);
\fill[formatamber] (3.290,-0.800) rectangle ++(.208,.27);
\fill[failgray] (3.525,-0.800) rectangle ++(.208,.27);
\fill[formatamber] (3.760,-0.800) rectangle ++(.208,.27);
\fill[passgreen] (3.995,-0.800) rectangle ++(.208,.27);
\fill[failgray] (4.230,-0.800) rectangle ++(.208,.27);
\fill[formatamber] (4.465,-0.800) rectangle ++(.208,.27);
\fill[formatamber] (4.700,-0.800) rectangle ++(.208,.27);
\fill[formatamber] (4.935,-0.800) rectangle ++(.208,.27);
\fill[formatamber] (5.170,-0.800) rectangle ++(.208,.27);
\fill[passgreen] (5.405,-0.800) rectangle ++(.208,.27);

\draw[black!45] (2.79,-.88)--(2.79,.31);
\fill[passgreen] (-.55,-1.18) rectangle ++(.16,.16);\node[anchor=west] at (-.34,-1.1) {Frozen task pass};
\fill[formatamber] (1.72,-1.18) rectangle ++(.16,.16);\node[anchor=west] at (1.93,-1.1) {Abstention + format failure};
\fill[failgray] (.65,-1.58) rectangle ++(.16,.16);\node[anchor=west] at (.86,-1.5) {Other automatic failure; semantic review pending};
\end{tikzpicture}
\caption{Paired outcomes retain the case identities hidden by equal totals.
Article order within each half follows the frozen cohort. Green is the original
automatic pass; amber is a post-hoc diagnostic for unanswerable cases beginning
with \texttt{unknown} but failing the contract. Gray does not mean proven factual
error: it includes lexical and citation mismatch. Standard/full each pass 10/24;
skip4 passes a different 10/24 and loses two standard-passing cases. No diagnostic
color changes the frozen verdict.}
\label{fig:research-quality}
\end{figure}


The two failures differ. Case 08 expands the correct ``October
1973'' into a sentence about the oil crisis, failing exact-reference scoring
without demonstrating a factual error. Case 16 gives ``pharma'' with the
placeholder citation \texttt{[source\_id:0]}, a real locator failure. The
interpretation must retain that distinction instead of describing both as
hallucinations or dismissing both as metric noise.

On the 12 unanswerable questions, standard produces ten responses beginning
with \texttt{unknown} but with extra/invalid citation text and two unsupported
substantive answers. All fail the strict contract. Ten are not automatically
ten hallucinations. The raw responses and review packet preserve the distinction.
The frozen counts are unchanged. Earlier ten-case development passes therefore
did not establish robust abstention or citation behavior on separate articles.

\subsection{RQ4 Findings: Publication, Answer Time and Prediction Error}
\textbf{The reservation remains after ingestion finishes.} At arrival zero,
E2 shared and 2+2 publication times are 14.2 and 7.5\,s for incremental updates,
and 50.2 and 26.8\,s for rebuilds. New answers finish at 288.5 and 413.7\,s for
incremental updates, and 304.6 and 420.8\,s for rebuilds. These medians cover
three different documents. Shared execution finishes the new answer earlier
in all 12 matched E2 pairs. Publication precedes current completion in all 24
concurrent episodes, but the fixed partition still limits the old and new
requests. Deferring ingestion protects the old request while delaying publication.
We do not test returning the reserved cores to inference after ingestion ends.

\textbf{The timing model has limited predictive value.} E3 correctly predicts
that shared execution is faster than 2+2 in six pairs. Mean absolute error is
2.54\,s for current completion, 0.50\,s for publication and 22.22\,s for new-answer
completion. The last value is a mean relative error of 6.49\% from update
arrival. A simple baseline using each policy's E2 median has an error of
21.48\,s. This baseline was added after E3 and was not part of the frozen
comparison. The phase model does not show better accuracy. In every E3 case,
publication finishes before the new request arrives and the occupied slot
limits progress. We therefore have not tested a boundary between regimes.

All 48 E2/E3 current-answer and retrieval checks pass. Of the new outputs,
32 pass the strict answer contract. These episodes cover six distinct
answerable questions. Two produce longer answers that fail exact matching
across their configurations. Human review is needed to assess their meaning.
The experiments do not test abstention. The earlier application control below uses a sequential arrival rule.

\subsubsection{Earlier Sequential Application Control}
\label{sec:joint}
\paragraph{Publication and answer times.}
All 12 application episodes complete. Current answers pass their supported
checks and retrieval locates the new source in v2 in every episode. Publication and retrieval therefore work in these episodes. We still need
the answer checks to determine whether the new outputs are valid.

\begin{table}[t]
\caption{Median times from update arrival across three matched episode blocks
with different new documents. $T_c$ is current completion, $T_p$ is publication and $T_f$ is new-answer completion. These are descriptive medians, not confidence bounds.}
\label{tab:joint}
\centering
\begin{tabular}{lrrr}
\toprule
Policy & $T_c$ & $T_p$ & $T_f$\\
\midrule
Standard concurrent & 76.2 & 14.7 & 189.8 \\
Standard deferred & 70.7 & 76.0 & 189.0 \\
Full prototype & 78.5 & 10.3 & 198.9 \\
Skip4 prototype & 76.6 & 9.9 & 189.6 \\
\bottomrule

\end{tabular}
\end{table}
The full prototype publishes earlier than standard concurrent (10.3 versus
14.7\,s median), but its current and fresh answers finish later (78.5 versus
76.2\,s for current answers and 198.9 versus 189.8\,s for new answers). Skip4 publishes at 9.9\,s, with fresh completion
189.6\,s, essentially the same scale as standard. It is already ineligible under
the separate-document paired gate. Deferral protects current service (70.7\,s)
but delays publication to 76.0\,s. None satisfies the frozen automatic joint
benefit criterion.

Figure~\ref{fig:research-episode} shows the timing of one measured episode. In every
concurrent episode, publication precedes current-answer completion. Under the
frozen single-slot trace, Equation~\ref{eq:critical-path} is then governed by
current inference and the next request rather than publication. The measured result is that publication is already complete while
current inference remains active, so accelerating publication does not advance
the next request in these episodes. It does not
prove that publication speed is irrelevant for earlier-arriving independent
requests, multiple slots or another workload.
\begin{figure*}[t]
\centering
\begin{tikzpicture}[font=\scriptsize,>=Latex]
\fill[answerblue] (0,.72) rectangle ++(.25,.16);\node[anchor=west] at (.33,.80) {Current inference};
\fill[updatepurple] (3.1,.72) rectangle ++(.25,.16);\node[anchor=west] at (3.43,.80) {Incremental publication};
\fill[freshorange] (7.0,.72) rectangle ++(.25,.16);\node[anchor=west] at (7.33,.80) {New-version inference};
\node[anchor=east,font=\scriptsize] at (-.18,-0.120) {Standard concurrent};
\fill[answerblue] (0,0.000) rectangle (4.1315,0.200);
\fill[updatepurple] (0.0000,-0.250) rectangle (0.7919,-0.050);
\draw[black!45,densely dotted] (4.1315,0.100) -- (4.1402,0.100);
\fill[freshorange] (4.1402,0.000) rectangle (10.2518,0.200);
\node[anchor=west,font=\scriptsize] at (10.3518,0.100) {189.8s};
\draw[updatepurple] (0.7919,-0.300) -- (0.7919,-0.460);
\node[anchor=north,font=\scriptsize,text=updatepurple] at (0.7919,-0.460) {14.7s};
\node[anchor=east,font=\scriptsize] at (-.18,-1.240) {Standard deferred};
\fill[answerblue] (0,-1.120) rectangle (3.8336,-0.920);
\fill[updatepurple] (3.8336,-1.370) rectangle (4.1127,-1.170);
\draw[black!45,densely dotted] (4.1127,-1.020) -- (4.1182,-1.020);
\fill[freshorange] (4.1182,-1.120) rectangle (10.2037,-0.920);
\node[anchor=west,font=\scriptsize] at (10.3037,-1.020) {189.0s};
\draw[updatepurple] (4.1127,-1.420) -- (4.1127,-1.580);
\node[anchor=north,font=\scriptsize,text=updatepurple] at (4.1127,-1.580) {76.2s};
\node[anchor=east,font=\scriptsize] at (-.18,-2.360) {Full prototype};
\fill[answerblue] (0,-2.240) rectangle (4.2989,-2.040);
\fill[updatepurple] (0.0000,-2.490) rectangle (0.5554,-2.290);
\draw[black!45,densely dotted] (4.2989,-2.140) -- (4.3068,-2.140);
\fill[freshorange] (4.3068,-2.240) rectangle (10.7381,-2.040);
\node[anchor=west,font=\scriptsize] at (10.8381,-2.140) {198.9s};
\draw[updatepurple] (0.5554,-2.540) -- (0.5554,-2.700);
\node[anchor=north,font=\scriptsize,text=updatepurple] at (0.5554,-2.700) {10.3s};
\node[anchor=west,font=\scriptsize,text=red!70!black] at (4.4068,-2.530) {correct fact; citation-format failure};
\node[anchor=east,font=\scriptsize] at (-.18,-3.480) {Skip4 prototype};
\fill[answerblue] (0,-3.360) rectangle (4.1341,-3.160);
\fill[updatepurple] (0.0000,-3.610) rectangle (0.5432,-3.410);
\draw[black!45,densely dotted] (4.1341,-3.260) -- (4.1429,-3.260);
\fill[freshorange] (4.1429,-3.360) rectangle (10.2359,-3.160);
\node[anchor=west,font=\scriptsize] at (10.3359,-3.260) {189.6s};
\draw[updatepurple] (0.5432,-3.660) -- (0.5432,-3.820);
\node[anchor=north,font=\scriptsize,text=updatepurple] at (0.5432,-3.820) {10.1s};

\draw[->,black!60] (0,-4.28)--(11.15,-4.28);
\foreach \t in {0,50,100,150,200}{\draw[black!60] (\t*.054,-4.22)--(\t*.054,-4.34);\node[anchor=north] at (\t*.054,-4.37) {\t};}
\node at (5.4,-4.92) {Elapsed seconds since update arrival};
\end{tikzpicture}
\caption{Measured critical path, first predeclared matched episode (Normandy),
not a timeline assembled from unrelated medians. Purple ticks label publication;
right labels mark new-answer completion. Publication precedes current-answer
completion for all concurrent arms, so accelerating publication does not release
the serial inference slot earlier. Deferral instead delays publication until the
current answer, leaving fresh completion similar in this episode. The full
prototype answers ``France'' but emits an invalid citation prefix; its timing
is preserved as a failed task. Dotted gaps are post-gate retrieval/submission.}
\label{fig:research-episode}
\end{figure*}


\paragraph{Answer checks.}
The first two added documents ask about Normandy's country and Southern
California's abbreviation. Standard answers ``France'' and ``SoCal'' with valid
locators. Full prototype's first answer has the correct fact but an extra
\texttt{source\_id:} prefix in the locator, which fails the contract. The third
asks what kind of economy Victoria has. All arms generate a relevant elaboration
containing the gold ``highly diversified'', but fail exact-reference scoring.
The longer answer fails the scoring rule. Human review is needed before
interpreting it as an error in the economic description.

The original all-episode quality gate fails for standard as well as the
candidates. We retain that failed result rather than relax the rule after reading the outputs. Separately, even a more permissive interpretation of that
shared elaboration would not make full prototype dominate the measured latency
vector: it publishes faster while answering later. Timing and answer validity are separate results. A new semantic assessment could change quality labels
but would not change these measured timestamps or establish statistical benefit.

The current response uses old pinned oracle evidence, whereas the new response
uses actual retrieval. We therefore claim a versioned-update prototype with one
retrieved fresh question per episode, not a fully validated production RAG
trace. The workload demonstrates neither domain-specific household deployment
utility nor a real-time service contract. Observed response times are in minutes.

